\documentclass[../../main.tex]{subfiles}
\begin{document}

\begin{flushright}
  \footnotesize\textit{【自動文字起こし・要確認】}
\end{flushright}

{\bf [解]} (1) 右のように回転軸 $l$ をとる．(図1)，(図2) のどちらのように回転させても，
\begin{align}
  V = (ac + \dfrac{1}{4} \pi (a^2 c)) b
\end{align}

\paragraph{(2)}

$\alpha = a+c$, $\beta = ac$ とおく．$a, c$ は0より大きい実数だから，
\begin{align}
  \begin{cases} \alpha^2 - 4\beta \ge 0 \\ 0 < \alpha \\ \beta > 0
  \end{cases} \cdots \text{①}
\end{align}
である．$a+b+c = 1$ から $b = 1 - \alpha$ で，$0 < b$ から
\begin{align}
  \alpha < 1 \cdots \text{②}
\end{align}
となる．$V$ に代入して
\begin{align}
  V &= (\beta + \dfrac{1}{4}\pi (\alpha^2 - 2\beta))(1-\alpha) \\
  &= (1-\alpha) \left\{ (1 - \dfrac{1}{2}\pi)\beta + \dfrac{1}{4}\pi \alpha^2 \right\} \cdots \text{③}
\end{align}
である．① $\land$ ②を図示すると右図斜線部で，③から $V$ は $\beta$ の単調減少関数だから，$\alpha$ を固定すると，
\begin{align}
  (1-\alpha) \left\{ (1-\dfrac{1}{2}\pi)\dfrac{1}{4}\alpha^2 + \dfrac{1}{4}\pi \alpha^2 \right\} \le V < \dfrac{1}{4}\pi \alpha^2 (1-\alpha) \cdots \text{④}
\end{align}
である．④の左辺 $f(\alpha)$，右辺 $g(\alpha)$ とおく．(境界は $\beta = \dfrac{1}{4}\alpha^2$ のみ含む)
\begin{align}
  \begin{cases} f(\alpha) = (1 + \dfrac{1}{8}\pi) \alpha^2 (1-\alpha) \\ g(\alpha) = \dfrac{1}{4}\pi \alpha^2 (1-\alpha)
  \end{cases} \cdots \text{⑤}
\end{align}
であり，$h(\alpha) = \alpha^2 (1-\alpha)$ とすると，$h'(\alpha) = 2\alpha - 3\alpha^2 = \alpha(2-3\alpha)$ から下表をうる．

\begin{table}[htb]
  \centering
  \begin{tabular}{|c|c|c|c|c|c|}
    \hline
    $\alpha$ & 0 & $\cdots$ & $2/3$ & $\cdots$ & 1 \\
    \hline
    $h'$ & & $+$ & 0 & $-$ & \\
    \hline
    $h$ & 0 & $\nearrow$ & $4/27$ & $\searrow$ & 0 \\
    \hline
  \end{tabular}
  \caption{$h(\alpha)=\alpha^2(1-\alpha)$の増減表}
\end{table}

したがって，$0 < h(\alpha) \le \dfrac{4}{27}$ だから ④，⑤とあわせて，
\begin{align}
  0 < V < \dfrac{1}{4}\pi \dfrac{4}{27} = \dfrac{\pi}{27}
\end{align}
\begin{align}
  \therefore 0 < V < \dfrac{\pi}{27}
\end{align}

\begin{figure}[htb]
  \centering
  \begin{tikzpicture}
    \draw[->] (-0.5,0) -- (2,0) node[right] {$\alpha$};
    \draw[->] (0,-0.5) -- (0,2) node[above] {$\beta$};
    \draw[domain=0:1, smooth, variable=\x] plot ({\x}, {\x*\x/4});
    \node at (1.5, 1) {$\beta = \frac{1}{4}\alpha^2$};
    \fill[pattern=north east lines] (0,0) -- plot[domain=0:1] ({\x}, {\x*\x/4}) -- (1,0) -- cycle;
  \end{tikzpicture}
  \caption{①$\land$②を満たす領域（斜線部，境界は$\beta=\dfrac{1}{4}\alpha^2$のみ含む）}
\end{figure}

\end{document}
